\section{Modalidades de Trabalho, Tecnologia e Segurança}

\subsection{Modelo de Trabalho}

Com a declaração da pandemia de Covid-19, a instituição implementou um plano de contingência para avaliar como o trabalho remoto em larga escala poderia ser utilizado na instituição. A adoçou do modelo foi realizada após rápida simulação, com a instituição funcionando neste modelo a partir do dia 23 de março do ano de 2020 \cite{bndes_ra2020}. Todavia, há de se observar que o órgão sempre apresentara em sua política interna preocupações com eventualidades semelhantes, contando com `` planejamento prévio necessário para eventos de continuidade de negócios e ferramentas de tecnologia da informação (TI) adequadas para permitir, de forma segura, o trabalho remoto de nossos empregados'' \cite{bndes_ra2020}. O \citeonline{bndes_ra2020} cita que o evento pandêmico fortaleceu as orientações aos colaboradores referentes assuntos de  procedimentos de acesso remoto e uso de ferramentas de reunião virtual, ergonomia no teletrabalho e procedimentos para segurança da informação.

A partir de agosto de 2022, a instituição implantou o \textbf{Programa Piloto de Trabalho Remoto Híbrido}, ``visando modernizar as práticas de gestão, ampliar o alinhamento à agenda ASG, aumentar o engajamento e a produtividade de (...) empregados e a qualidade das entregas para a sociedade, em linha com movimento global de mudanças nas relações de trabalho'' \cite{bndes_ra2022}. O programa piloto foi executado até abril de 2023, com um total de 2.152 empregados, dos quais a maior parcela executava suas funções remotamente duas vezes por semana \cite{bndes_ra2024}; em dezembro de 2025, este número 2.902 empregados ativos \cite{bndes_ra2025}.

Este modelo exige atenção aos tópicos de Segurança da Informação, ao que a instituição dá especial atenção, com campanhas recorrentes de educação em cibersegurança,  implantação de serviços especializados em segurança da informação para monitoramento de eventos, gestão de vulnerabilidades, tratamento de incidentes e inteligência de ameaças cibernéticas; \cite{bndes_ra2025}. 

Há de se observar que este tema tem sido de preocupação recorrente para a instituição, a qual coloca ``o aprimoramento contínuo dos
mecanismos para detecção, contenção e resposta a ameaças digitais tem sido prioridade em nosso planejamento estratégico de segurança da informação'' \cite{bndes_ra2025}. O BNDES realiza anualmente campanhas e ações voltadas ao aperfeiçoamento das suas capacidades de monitorar ameaças e reagir a incidentes e eventos que coloquem a segurança da informação da instituição em risco, contratando empresas especializadas em monitoramento, tratamento de incidentes e gestão de vulnerabilidades, realizando exercícios e testes de simulação de ataques para assegurar a eficácia de sua política de segurança e das medidas adotadas \cite{bndes_ra2022,bndes_ra2024,bndes_ra2025}.